\documentclass[10pt,a4paper]{amsart}
\usepackage[margin=19mm]{geometry}
\usepackage{amsmath,amssymb,mathtools}
\usepackage[scr=rsfs,cal=euler]{mathalfa}
\usepackage{xfrac}
\usepackage[unicode,hidelinks,bookmarks=false]{hyperref}
\newcommand{\ie}{i.e.\ }
\newcommand{\cf}{cf.\ }
\newcommand{\resp}{respectively\ }
\newcommand{\Subsection}{\subsection}
\providecommand{\nobreakdash}{\nobreak\hbox{-}}
\setlength{\emergencystretch}{2em}
\multlinegap0pt
\usepackage{amssymb}
\usepackage{tikz}
\usepackage{float}
\usepackage[shortlabels]{enumitem}
\newtheorem{theorem}{Theorem}
\newcommand*{\bbbxi}{\boldsymbol{\xi}}
\title{Reducible Iterated Graph Systems: \\ multiscale-freeness and multifractals}
\author{Nero Ziyu Li, Frank Xin Hu, Thomas Britz}
\date{Source: arXiv:2506.18073v2 (CC BY 4.0)}
\begin{document}
\maketitle

\noindent\emph{Source and licence.} Reducible Iterated Graph Systems: \\ multiscale-freeness and multifractals. Version arXiv:2506.18073v2 (CC BY 4.0), distributed by arXiv under the Creative Commons Attribution 4.0 International licence, http://creativecommons.org/licenses/by/4.0/, as recorded in the arXiv licence metadata for this exact version and shown on the abstract page https://arxiv.org/abs/2506.18073v2. Full licence notice: \texttt{licenses/arxiv-2506.18073-CC-BY-4.0.txt}.\par\smallskip

\noindent\textit{Editorial scope.} Counted unit: the article's theorem theorem:distance (source0.tex lines 1596--1615). The article's complete original proof of it is delivered with it (source0.tex lines 1617--1735, one \texttt{proof} environment of 119 lines). No mathematics is written, repaired or re-derived: the statement and the proof are the article's own lines, in the article's own notation, and no retained line is substituted or re-flowed.  No further source-local context is needed: every cross-reference of every retained line resolves inside the two delivered spans.  Nothing was omitted from any retained span, so the package carries no omission record.  No retained line contains a citation, so the wrapper carries no bibliography and no bibliography entry.  No retained line contains a figure environment, an image-inclusion command or drawing code, so no image is delivered and the package declares no asset dependency.  Editorial formatting only, in addition: an AMS \texttt{amsart} preamble that restates the article's own declarations and macros used by the retained lines, namely the environments \texttt{theorem} on separate counters, together with the standard packages those lines need.  The retained environments print in this document as Theorem 1 in order of appearance, whereas the article prints the same items with the numbering of its own full document.  The complete native source of the article as distributed in the arXiv e-print for this exact version is bundled unchanged as \texttt{sources/180296/source0.tex}.
\begin{theorem}\label{theorem:distance}
For every colour $j\in\mathfrak r_{\mathbf M}(\iota)$ there exists a constant $c_j\in\mathbb R_+$ such that
\[
\frac{1}{\Lambda_{\mathcal D}(j)^n}
\min_{\mathbf D_1,\ldots,\mathbf D_n\in\mathcal D}
\bigg\|
\bbbxi_j\prod_{m=1}^n \mathbf D_m
\bigg\|_1
\to c_j
\]
as $n\to\infty$.
Moreover, for every sequence $\mathbf D_1,\ldots,\mathbf D_n$ with $\mathbf D_m\in\mathcal D$ for each $m\in[n]$,
\[
\bigg\|
\bbbxi_j\prod_{m=1}^n \mathbf D_m
\bigg\|_1
\succeq
\Lambda_{\mathcal D}(j)^n.
\]
\end{theorem}

\begin{proof}
Fix $j\in\mathfrak r_{\mathbf M}(\iota)$ and write
$A(\mathbf D):=A_j(\mathbf D)$ and $\lambda:=\Lambda_{\mathcal D}(j)$.

We first show that $R_j$ is forward invariant under every $\mathbf D\in\mathcal D$.
If $a\in R_j$ and $[\mathbf D]_{ab}>0$, choose $\mathbf E\in\mathcal D$ and a path
$j=i_0,i_1,\ldots,i_m=a$ such that $[\mathbf E]_{i_{r-1}i_r}>0$ for all $r\in[m]$.
Replacing the $a$-th row of $\mathbf E$ by the $a$-th row of $\mathbf D$ preserves the path and adds the edge $a\to b$.
Hence $b\in R_j$.

Therefore, for every sequence $\mathbf D_1,\ldots,\mathbf D_n\in\mathcal D$,
\[
\bigg\|
\bbbxi_j\prod_{m=1}^n \mathbf D_m
\bigg\|_1
=
\bigg\|
\bbbxi_j\prod_{m=1}^n A(\mathbf D_m)
\bigg\|_1.
\]

Let $A^*:=A(\mathbf D_j^*)$ and let $\mathbf v_j^*>0$ be its right Perron vector, so
$A^*\mathbf v_j^*=\lambda\mathbf v_j^*$.
We claim that $A(\mathbf D)\mathbf v_j^*\ge \lambda\mathbf v_j^*$ for every $\mathbf D\in\mathcal D$.
Assume not.
Then for some $\mathbf D\in\mathcal D$ and some $a_0\in R_j$ we have
$[A(\mathbf D)\mathbf v_j^*]_{a_0}<\lambda[\mathbf v_j^*]_{a_0}$.
Let $\mathbf D'$ be obtained from $\mathbf D_j^*$ by replacing its $a_0$-th row by the $a_0$-th row of $\mathbf D$.
By assumption, $A(\mathbf D')$ is primitive.
Also $A(\mathbf D')\mathbf v_j^*\le \lambda\mathbf v_j^*$, and the inequality is strict in the coordinate $a_0$.
Choose $q$ with $A(\mathbf D')^q>\mathbf 0$ and set $\mathbf u:=A(\mathbf D')^q\mathbf v_j^*$.
Then $\mathbf u>\mathbf 0$ and $A(\mathbf D')\mathbf u<\lambda\mathbf u$ coordinatewise.
The Collatz--Wielandt formula gives $\rho(A(\mathbf D'))<\lambda$.
Since $A(\mathbf D')$ is primitive, its whole trim is reachable from $j$, so
\[
\max_{\ell\in\mathfrak R_{\mathbf D'}(j)}
\rho\big(\mathbf B_\ell(\mathbf D')\big)
=
\rho\big(A(\mathbf D')\big)
<
\lambda,
\]
contrary to the definition of $\lambda$.
This proves the claim.

Iterating the claim gives
$\prod_{m=1}^n A(\mathbf D_m)\mathbf v_j^*\ge \lambda^n\mathbf v_j^*$.
Hence
\[
\bbbxi_j\prod_{m=1}^n A(\mathbf D_m)\mathbf v_j^*
\ge
\lambda^n[\mathbf v_j^*]_j.
\]
Since $\|\mathbf w\|_1\ge (\mathbf w\mathbf v_j^*)/\|\mathbf v_j^*\|_\infty$ for every non-negative row vector $\mathbf w$, we obtain
\[
\bigg\|
\bbbxi_j\prod_{m=1}^n \mathbf D_m
\bigg\|_1
\ge
\frac{[\mathbf v_j^*]_j}{\|\mathbf v_j^*\|_\infty}\lambda^n.
\]
This proves the lower bound.

For the limit, define the Bellman operator
$\mathcal F_j:\mathbb R_{\ge 0}^{|R_j|}\to\mathbb R_{\ge 0}^{|R_j|}$ by
$[\mathcal F_j(\mathbf x)]_a:=\min_{\mathbf D\in\mathcal D}[A(\mathbf D)\mathbf x]_a$.
Because the rows are chosen independently,
\[
[\mathcal F_j^n(\mathbf 1)]_j
=
\min_{\mathbf D_1,\ldots,\mathbf D_n\in\mathcal D}
\bigg\|
\bbbxi_j\prod_{m=1}^n \mathbf D_m
\bigg\|_1.
\]
The claim above implies $\mathcal F_j(\mathbf v_j^*)\ge \lambda\mathbf v_j^*$.
Since $\mathbf D_j^*\in\mathcal D$, we also have
$\mathcal F_j(\mathbf v_j^*)\le A^*\mathbf v_j^*=\lambda\mathbf v_j^*$.
Thus $\mathcal F_j(\mathbf v_j^*)=\lambda\mathbf v_j^*$.

Set $\mathcal G_j:=\lambda^{-1}\mathcal F_j$.
Choose $a_0,b_0>0$ so that $a_0\mathbf v_j^*\le \mathbf 1\le b_0\mathbf v_j^*$.
Then $a_0\mathbf v_j^*\le \mathcal G_j^n(\mathbf 1)\le b_0\mathbf v_j^*$ for all $n$.
Write $\mathbf y_n:=\mathcal G_j^n(\mathbf 1)$ and
$\beta_n:=\max_a[\mathbf y_n]_a/[\mathbf v_j^*]_a$.
Then $(\beta_n)$ is decreasing, hence converges.

Let $\mathbf y$ be an accumulation point of $(\mathbf y_n)$.
If $\mathbf y\ne \beta\mathbf v_j^*$, where $\beta:=\lim\beta_n$, then
$\mathbf y<\beta\mathbf v_j^*$ in at least one coordinate.
Choose $q$ with $(A^*)^q>\mathbf 0$.
Since $\mathcal F_j(\mathbf x)\le A^*\mathbf x$ for every $\mathbf x\ge \mathbf 0$, we get
\[
\mathcal G_j^q(\mathbf y)
\le
\lambda^{-q}(A^*)^q\mathbf y
<
\beta\lambda^{-q}(A^*)^q\mathbf v_j^*
=
\beta\mathbf v_j^*,
\]
which contradicts the definition of $\beta$.
Hence every accumulation point of $(\mathbf y_n)$ equals $\beta\mathbf v_j^*$, so
$\mathbf y_n\to \beta\mathbf v_j^*$.

Therefore
\[
\frac{1}{\lambda^n}
\min_{\mathbf D_1,\ldots,\mathbf D_n\in\mathcal D}
\bigg\|
\bbbxi_j\prod_{m=1}^n \mathbf D_m
\bigg\|_1
=
[\mathbf y_n]_j
\to
\beta[\mathbf v_j^*]_j.
\]
Since $\mathbf v_j^*>0$, the limit is strictly positive.
\end{proof}

\end{document}
